% Generated draft: requires mathematical and source review.
\documentclass[11pt]{article}
\usepackage[margin=25mm]{geometry}
\usepackage{fontspec}
\setmainfont{Latin Modern Roman}
\usepackage{amsmath,amssymb,mathtools,unicode-math}
\setmathfont{Latin Modern Math}
\usepackage{enumitem,needspace,xurl,hyperref}
\setlength{\parindent}{0pt}
\setlength{\parskip}{5pt}
\setlength{\emergencystretch}{3em}
\newcommand{\N}{\mathbb{N}}
\newcommand{\Z}{\mathbb{Z}}
\newcommand{\Q}{\mathbb{Q}}
\newcommand{\R}{\mathbb{R}}
\newcommand{\C}{\mathbb{C}}
\newcommand{\F}{\mathbb{F}}
\newcommand{\A}{\mathbb{A}}
\newcommand{\PP}{\mathbb{P}}
\newcommand{\E}{\mathbb{E}}
\newcommand{\eps}{\varepsilon}
\newcommand{\dd}{\,\mathrm{d}}
\begin{document}
\section*{UnsolvedMath ID: 20001927}
\textbf{Source problem code:} catalogue entry
\subsection{Definitions and mathematical statement}
Let \(M\) be a smooth compact manifold without boundary, of dimension \(n\geq2\), and let \(g\) be a Lorentz metric. ``Globally conformal to constant curvature'' means that, after a global identification if necessary,
\[
 g=e^{2u}g_0,
\]
where \(u\in C^\infty(M,\R)\) and \(g_0\) is a Lorentz metric of constant sectional curvature. Question \(9.3.2\), attributed to [Sanchez2013], asks whether every such \(g\) is geodesically complete: every maximal affinely parametrized geodesic, of any causal type, is defined on all of \(\R\). Compactness alone does not imply Lorentzian geodesic completeness.

The supplied aid reports these partial conclusions:
\begin{enumerate}
\item \(g\) is null-geodesically complete in every dimension. Along a null geodesic the conformal affine-parameter change is proportional to \(e^{2u}\), which is bounded above and below by positive constants on \(M\); this uses the completeness of the compact constant-curvature representative in the applicable Lorentz setting.
\item In dimension \(2\), \(g\) is fully geodesically complete.
\item In any dimension, \(g\) is fully complete if a finite cover of \((M,g_0)\) admits an everywhere timelike Killing vector field. This hypothesis concerns the representative \(g_0\), not a Killing field for every conformal metric. On the cover it becomes a timelike conformal Killing field for \(g\).
\item If \(g_0\) is flat and its linear holonomy
\[
 \rho:\pi_1(M)\to\operatorname{O}(1,n-1)
\]
has compact closure, then \(g\) is fully geodesically complete. The aid uses the resulting timelike parallel direction after a suitable finite cover.
\end{enumerate}
Consequently, any geodesically incomplete example with a flat representative must have linear holonomy with noncompact closure, equivalently an unbounded linear-holonomy group in a fixed linear representation. Null completeness by itself leaves timelike and spacelike completeness to be addressed, and compact holonomy is a sufficient condition rather than a necessary condition for completeness. The aid reports the general higher-dimensional problem apparently open in the literature checked, without a complete search bibliography.

The original extraction also includes contextual commentary that noncompact homogeneous indefinite-signature manifolds need not be complete, referring to [Sanchez2013, Example \(2\), Section \(4\)], and suggests asking which algebraic properties of the isometry group might imply completeness. The verified record boundary identifies this as commentary, not a hypothesis or answer to Question \(9.3.2\). The subsequent Section \(10\) and \(10.1\) headings on integrability and ergodicity of higher-genus surface geodesic flows are further extraction spillover.
\subsection{Short English statement}
The question asks whether a compact Lorentz metric globally conformal to a constant-curvature metric must be complete. The aid proves null completeness, the full surface case, and completeness under a finite-cover timelike-Killing condition. In the flat case, an incomplete example would require unbounded linear holonomy; the general higher-dimensional question remains open in the supplied report.
\subsection{Sources}
\begin{itemize}
\item Local UnsolvedMath registry, permanent ID \(20001927\), source code \texttt{catalogue entry}; canonical collection \texttt{aim-geometry-notes.json}, notes[264]. Source item \(9.3\), actual question \(9.3.2\), attributed to [Sanchez2013]. Statement status \texttt{corrected\_verified}.
\item AIM workshop ``Geodesics,'' Burns--Matveev report ``Open Problems and Questions about Geodesics'': \url{https://aimath.org/pastworkshops/geodesicsproblems.pdf}. Supplied verification identifies official page \(23\), PDF index \(22\), and excludes the following homogeneous-metric commentary and Section \(10\) headings from the mathematical question. The PDF was not opened in this session.
\item [Sanchez2013], including Example \(2\) in Section \(4\) for noncompact homogeneous indefinite metrics; no complete title or publication data are supplied. The original literature and remarks fields are empty; no proposer-field value or year is supplied.
\item Supplied research attempt \(1\), \texttt{partial\_result}, run status ``open,'' rigor label ``proved,'' novelty confidence ``low.'' Registry status \texttt{partially\_solved}, classification \texttt{PARTIAL-PROGRESS}, updated 14 August 2026. Full authoritative references for the compact constant-curvature, surface and conformal-Killing completeness inputs are absent.
\end{itemize}
\end{document}
